All nine are \textbf{off-the-shelf deep-learning models for time series}, taken from the Python
library \textbf{Darts} (\texttt{unit8co/darts}). The architectures are the library's, unmodified ---
each is a published model used as Darts implements it. What we supply is the data, the
configuration, and for six of the nine a loss function of our own (Section~\ref{sec:losses}).

They are \textit{global} models in the forecasting sense: a single set of parameters is pooled
across all 64,818 cells. The closest familiar analogue is a pooled panel estimator --- a cell with
almost no conflict history is still forecast using structure learned from the whole grid. Each
cell's own level and scale are restored separately, by the normalisation described in
Section~\ref{sec:architectures}.

\begingroup
\catcode`\_=12  % underscores are ordinary characters in this table
\begin{table}[h!]
\centering
\small
\renewcommand{\arraystretch}{1.2}
\caption{The nine delivered models.}
\label{tab:models}
\begin{tabular}{@{}p{2.9cm}p{1.7cm}p{1.9cm}p{2.7cm}p{3.1cm}p{1.5cm}@{}}
\toprule
\textbf{Model} & \textbf{Type} & \textbf{Loss} & \textbf{Predictors}
  & \textbf{Hyperparameters} & \textbf{Features} \\ \midrule
dark_river       & N-BEATS & MSE       & Outcomes only   & \cfg{dark_river}       & \qs{dark_river} \\
old_rules        & N-BEATS & Spotlight & Outcomes only   & \cfg{old_rules}        & \qs{old_rules} \\
blue_ocean       & N-BEATS & Spotlight & + 69 covariates & \cfg{blue_ocean}       & \qs{blue_ocean} \\ \addlinespace
silent_fox       & TSMixer & MSE       & Outcomes only   & \cfg{silent_fox}       & \qs{silent_fox} \\
dancing_monkey   & TSMixer & Spotlight & Outcomes only   & \cfg{dancing_monkey}   & \qs{dancing_monkey} \\
brave_heart      & TSMixer & Spotlight & + 69 covariates & \cfg{brave_heart}      & \qs{brave_heart} \\ \addlinespace
red_hawk         & TiDE    & MSE       & Outcomes only   & \cfg{red_hawk}         & \qs{red_hawk} \\
dark_necessities & TiDE    & Spotlight & Outcomes only   & \cfg{dark_necessities} & \qs{dark_necessities} \\
golden_eagle     & TiDE    & Spotlight & + 69 covariates & \cfg{golden_eagle}     & \qs{golden_eagle} \\
\bottomrule
\end{tabular}
\end{table}
\endgroup

\noindent The link columns open each model's configuration files on GitHub, fixed at the delivery
commit on the title page: \textbf{Hyperparameters} is \texttt{config\_hyperparameters.py}, which sets the
architecture, the loss and the training schedule; \textbf{Features} is
\texttt{config\_queryset.py}, which declares the predictors the model was given.

The names are arbitrary labels, assigned according to the \texttt{adjective\_noun} convention
used across views-models, and carry no meaning.

\textbf{Within each block of three, exactly two settings vary} --- the loss function, and whether
the model sees covariates. Everything else is identical: layer widths, depth, training epochs,
learning rate, early stopping. So inside a block, each row differs from the one above it in a
single setting. Across blocks the architecture changes and the comparison is a different one.

\begin{itemize}
    \item \textbf{Outcomes only} --- the model's inputs are the past values of the three outcomes
          and nothing else. Purely autoregressive. Six models: the first two rows of each block.
    \item \textbf{+ 69 covariates} --- the same three histories plus 69 predictors: ACLED event and
          fatality counts, land cover shares, population and built-up area, infant and child
          mortality, V-Dem governance indices, SHDI development indices, travel time to cities,
          mineral and petroleum deposits, growing season and harvested area, and terrain. Three
          models: the third row of each block.
\end{itemize}

\subsection*{How the Covariates Are Conditioned}

Three things happen to a covariate between the data store and the model. Only the last is ours.

\textbf{Spatial handling.} Nothing arrives natively on the PRIO-GRID. Population and built-up
surface are aggregated by summing $60\times60$ source raster pixels per cell; V-Dem and SHDI are
attached to cells by polygon-to-centroid spatial joins. The source documentation describes these as
opinionated decisions rather than mechanical transforms, and they are not reversible from the
delivered data.

\textbf{Temporal handling.} There is no feature engineering --- no lags, decay functions, ratios or
derived indices; the datafactory places cross-feature derivation in the model layer by design
(ADR-014 §6). But the series are not left as-is either. Annually observed quantities are extended
across their twelve months, by one of two per-source rules (ADR-014 §5): V-Dem and SHDI hold the
annual value flat; population and built-up surface interpolate linearly between five-yearly epochs.
Outside a source's observed range, values are \emph{manufactured}: zero before the first
observation, and the last observation held constant after it (ADR-052). Where a source's coverage
is narrower than the grid, the gap is filled with zeros rather than nulls, deliberately, because
consumers assume float arrays without missing-value handling (ADR-047). SHDI is the one exception:
cells its source never covered keep a null (ADR-042).

\textbf{In the model configuration.} All 72 columns --- the three outcomes and all 69 covariates
--- then pass through an \texttt{asinh} transform before training, declared in each model's
hyperparameter configuration (linked in Table~\ref{tab:models}). The transform is inverted for the
outcomes, which is why the delivered forecasts are in counts.

\subsection*{A Consequence Worth Knowing Before Using the Covariate Models}

The zero-fill rule above has a specific effect on this delivery, and it is not visible in the data:
\textbf{a manufactured zero is indistinguishable from an observed zero.}

The store publishes the real observation window per source. Checked against this delivery's
partition:

\begin{table}[h!]
\centering
\small
\renewcommand{\arraystretch}{1.2}
\caption{Source coverage against the calibration partition. Months outside a source's observed
window are manufactured, not measured.}
\label{tab:coverage}
\begin{tabular}{@{}p{3.4cm}p{3.6cm}p{3.6cm}p{3.4cm}@{}}
\toprule
\textbf{Source} & \textbf{Observed from} & \textbf{In training (1990--2017)} & \textbf{In test (2018--2021)} \\ \midrule
ACLED (7 covariates) & January 2020 & \textbf{none --- all zeros} & real from January 2020 \\
V-Dem                & January 1989 & real & real \\
SHDI                 & January 1989 & real & real \\
Population           & January 1989 & real & real \\
Built-up surface     & January 1989 & real & real \\
\bottomrule
\end{tabular}
\end{table}

\noindent \textbf{The seven ACLED covariates were identically zero throughout training.} They carry
no information for any month the three covariate-using models were fitted on, and then become
non-zero partway through the forecast window. A feature that is constant in training and varies in
test is not merely uninformative --- the model has learned that it is always zero.

This affects \texttt{blue\_ocean}, \texttt{brave\_heart} and \texttt{golden\_eagle} only. The six
outcomes-only models do not use covariates at all. The other four sources cover the whole partition.

\textbf{What this note does not address.} Whether any of this conditioning is appropriate --- the
fill rules, the zero-for-missing convention, the spatial joins, the \texttt{asinh} transform, or
the treatment of annual variables at monthly resolution --- is outside the scope of a delivery
note. Nothing here endorses those choices and nothing here has examined them. They are recorded so
that someone evaluating these models knows what the models were given. Two limitations of the
record itself are worth stating: provenance is held per \emph{source}, not per covariate, so
reconstructing what happened to any one of the 69 means assembling it from the source mapping, the
relevant ADR and the run ledger; and the population and built-up-surface series include projected
rather than measured epochs at recent dates, delivered identically to measured ones --- beyond this
partition's end, but worth knowing if the window is ever extended.

Note that the two splits are not the same split. Six models are outcomes-only and six use the
Spotlight loss, but they are different sixes: \texttt{blue\_ocean}, \texttt{brave\_heart} and
\texttt{golden\_eagle} have covariates \textit{and} Spotlight, while \texttt{dark\_river},
\texttt{silent\_fox} and \texttt{red\_hawk} have neither.

\textbf{There is no model combining MSE with the covariates.} Three configurations exist, not four.
